I cannot honestly give the requested proof, because the lemma is **false as stated**.

Here is the rigorous resolution.

**Claim.** Under the hypotheses written in the prompt, the conclusion \(m_k=0\) does not follow.

**Proof.** For **every** polygon with vertices \(s_1,\dots,s_n\) and edge vectors \(a_i=s_{i+1}-s_i\) (with \(s_{n+1}=s_1\)), one has the telescoping identity
\[
\sum_{i=1}^n a_i
=
\sum_{i=1}^n (s_{i+1}-s_i)
=
0.
\]
Hence
\[
(m_1,\dots,m_n)=(1,\dots,1)
\]
is always a nonnegative integer solution of
\[
\sum_{i=1}^n m_i a_i=0.
\]

Therefore, if the proposed lemma were true, then **no** convex polygon could have an edge \(a_k\) with no parallel counterpart, because the universal choice \(m_i=1\) would force
\[
m_k=1=0,
\]
a contradiction.

But such convex polygons certainly exist; the simplest example is a triangle.

Take
\[
s_1=(0,0),\qquad s_2=(1,0),\qquad s_3=(0,1).
\]
Then \(P=\operatorname{conv}\{s_1,s_2,s_3\}\) is a convex polygon, and its counterclockwise edge vectors are
\[
a_1=s_2-s_1=(1,0),\qquad
a_2=s_3-s_2=(-1,1),\qquad
a_3=s_1-s_3=(0,-1).
\]
No two of \(a_1,a_2,a_3\) are parallel. In particular, for each \(k\), the vector \(a_k\) has no parallel counterpart of either sign among the other edges.

Now choose
\[
m_1=m_2=m_3=1.
\]
Then
\[
m_1a_1+m_2a_2+m_3a_3
=
(1,0)+(-1,1)+(0,-1)
=
(0,0).
\]
So all the hypotheses of the stated lemma hold, but
\[
m_1=m_2=m_3=1>0.
\]
This contradicts the claimed conclusion \(m_k=0\).

Thus the lemma is false. \(\square\)

---

**Where the proposed projection argument fails.**

The problematic step is the assertion that cancellation of the \( \hat a_k \)-component requires an antiparallel edge. That is not true.

In the triangle above,
\[
-a_1=(-1,0)=(-1,1)+(0,-1)=a_2+a_3.
\]
So \(a_1\) is canceled by **two nonparallel vectors**, not by a single vector parallel to \(-a_1\).

If we project onto
\[
\hat a_1=(1,0),\qquad \hat a_1^\perp=(0,1),
\]
then
\[
a_2\cdot \hat a_1=-1,\qquad a_3\cdot \hat a_1=0,
\]
so the \( \hat a_1 \)-component of \(a_1\) is canceled without any antiparallel copy of \(a_1\). And
\[
a_2\cdot \hat a_1^\perp=1,\qquad a_3\cdot \hat a_1^\perp=-1,
\]
so the perpendicular projection equation is simply a cancellation of one positive scalar by one negative scalar. The fact that the coefficients \(m_j\) are nonnegative does **not** imply that each projected term must vanish separately.

---

**What projection does prove.**

A correct statement is the following:

**Lemma.** If
\[
\sum_i m_i a_i=0,\qquad m_i\ge 0,
\]
and for some nonzero vector \(u\) there exists an index \(i\) with
\[
m_i>0
\quad\text{and}\quad
u\cdot a_i>0,
\]
then there exists an index \(j\) with
\[
m_j>0
\quad\text{and}\quad
u\cdot a_j<0.
\]

**Proof.** Otherwise \(u\cdot a_j\ge 0\) for every active index \(j\), and at least one active index has \(u\cdot a_i>0\). Then
\[
0
=
u\cdot \sum_i m_i a_i
=
\sum_i m_i (u\cdot a_i)
>
0,
\]
a contradiction. \(\square\)

So projection only shows that the active edge vectors cannot all lie in a single open half-plane through the origin. It does **not** force cancellation to occur by parallel or antiparallel direction classes.

---

If you want a version with **no accidental zero projection** at all, one can also use the five unit vectors at arguments \(0,\frac{2\pi}{5},\frac{4\pi}{5},\frac{6\pi}{5},\frac{8\pi}{5}\): their sum is \(0\) by the fifth-roots-of-unity identity, and their partial sums form a regular pentagon, hence a convex polygon. That gives a counterexample even when every other edge is neither parallel nor perpendicular to \(a_k\). ([dlmf.nist.gov](https://dlmf.nist.gov/1.11.iv))

So the correct conclusion is:

\[
\textit{the closure equation } \sum_i m_i a_i=0 \textit{ with } m_i\ge 0
\textit{ does not imply } m_k=0 \textit{ from “no parallel counterpart” alone.}
\]

If you want, I can help formulate a **corrected lemma** that uses additional hypotheses from \( \mathrm{SP1\text{--}4} \) and *is* actually provable.